\section{Statistical Decision Theory and Sufficiency}

\subsection{Fundamental Definitions}

\begin{definition}
    Let
\[
X:(\Omega,\mathcal{F})\longrightarrow(\mathcal{X},\mathcal{B})
\]
be the random observation. Its distribution is a probability measure on
the sample space $\mathcal{X}$; the realized data are a particular value
$x$. We use $P_\theta$ for the law of the whole observation when
describing a model indexed by an unknown parameter $\theta$.

For an i.i.d.\ sample, write
\[
X_1,\ldots,X_n\sim Q_\theta
\quad\text{independently},
\]
so the full sample has law
\[
P_\theta=Q_\theta^{\otimes n}.
\]
Here $Q_\theta$ is the \textbf{population distribution} of one unit, and
$n$ is the \textbf{sample size}. If each unit has $d$ recorded features,
then
\[
x=(x_1,\ldots,x_n)^\top\in\mathbb{R}^{n\times d}
\]
is a data matrix.

An observation may be a person, a measurement, or a more complicated
object such as an entire graph. The statistical goal is to learn about
the population, or a feature of it, from the observed sample.
\end{definition}

\begin{definition}[Parametric model]
A \emph{parametric model} is a family
\[
\mathcal{P}
=
\left\{
P_\theta:\theta\in\Theta\subseteq\mathbb{R}^k
\right\}
\]
with a fixed finite-dimensional parameter $\theta$. The set $\Theta$ is
the \emph{parameter space}. Inference may concern the whole parameter or
a specified function $\psi(\theta)$ of it.
\end{definition}

\begin{definition}[Dominated model]
A family is \emph{dominated} by a common sigma-finite measure $\nu$ if
\[
P_\theta\ll\nu
\qquad\text{for every }\theta.
\]
Then its densities are
\[
p_\theta=\frac{\mathrm{d}P_\theta}{\mathrm{d}\nu}.
\]
Lebesgue measure gives probability density functions; counting measure
gives probability mass functions. The reference measure is fixed across
the family.
\end{definition}

\begin{definition}[Identifiability]
A parametrization
\[
\{P_\theta:\theta\in\Theta\}
\]
is \emph{identifiable} if
\[
P_{\theta_1}=P_{\theta_2}
\quad\Longrightarrow\quad
\theta_1=\theta_2.
\]

Even in a non-identifiable parametrization, a target $\psi(\theta)$ can
be identifiable: it must have the same value whenever two parameters
give the same observation law.
\end{definition}

\begin{example}[Multinomial/Softmax regression]
A categorical response $Y\in \{1,\ldots,K\}$ and observed covariates
$x\in\mathbb{R}^d$, let
\[
p_k(x)
=
P(Y=k\mid x)
=
\frac{\exp(\beta_k^\top x)}
{\sum_{j=1}^K\exp(\beta_j^\top x)}.
\]
Adding the same vector $c$ to every $\beta_k$ multiplies the numerator and
denominator by $\exp(c^\top x)$, so it leaves all probabilities unchanged.
Choose the last class as the reference, and set
$\gamma_k=\beta_k-\beta_K$ for $k<K$ and $\gamma_K=0$. Then
\[
p_k(x)
=
\frac{\exp(\gamma_k^\top x)}
{1+\sum_{j=1}^{K-1}\exp(\gamma_j^\top x)},
\qquad k<K,
\]
and
\[
p_K(x)
=
\frac{1}
{1+\sum_{j=1}^{K-1}\exp(\gamma_j^\top x)}.
\]
This removes the common-shift ambiguity. Identifying the coefficient vectors
also requires enough variation in the design: on observed design points
spanning $\mathbb{R}^d$, the log odds
\[
\log\left(\frac{p_k(x)}{p_K(x)}\right)=\gamma_k^\top x
\]
determine $\gamma_k$.
\end{example}

\begin{definition}[Statistic]
A \emph{statistic} is a measurable function $T(X)$ of the observed data.
The rule
\[
T:\mathcal{X}\longrightarrow\mathcal{T}
\]
must be specified without using an unknown parameter. Its distribution
may nevertheless depend on that parameter. For real-vector data and
summaries, measurability is with respect to their Borel
sigma-algebras.
\end{definition}

\begin{definition}
    [Order Statistic]The order statistics
\[
X_{(1)}\leq\cdots\leq X_{(n)}
\]
are the sample values sorted in increasing order, retaining repeats if
there are ties. Suppose the observations are i.i.d.\ with CDF $F$ and
Lebesgue density $f$. For $1\leq i\leq n$, the $i$th order statistic has
density
\begin{equation}
g_i(x)
=
\frac{n!}{(i-1)!(n-i)!}
F(x)^{i-1}\bigl\{1-F(x)\bigr\}^{n-i}f(x),
\qquad\text{for Lebesgue-a.e.\ }x.
\end{equation}
\end{definition}

\begin{definition}[Exponential family]
A dominated family is an \emph{exponential family} if its densities have
the form
\begin{equation}
p_\theta(x)
=
\exp\left\{
\eta(\theta)^\top T(x)-\xi(\theta)
\right\}h(x),
\qquad
\xi(\theta)=A\bigl(\eta(\theta)\bigr),
\tag{5}
\end{equation}
where $T(x)\in\mathbb{R}^m$ and $h(x)\geq0$ do not depend on $\theta$,
and
\[
A(\eta)
=
\log\int
\exp\left\{\eta^\top T(x)\right\}
h(x)\,\mathrm{d}\nu(x).
\]

The \emph{natural parameter} is $\eta$. Its natural parameter space
$\mathcal{H}$ consists of those $\eta$ for which this integral is
positive and finite. In canonical form,
\[
p_\eta(x)
=
\exp\left\{
\eta^\top T(x)-A(\eta)
\right\}h(x).
\]
The function $A$ is the \emph{log-normalizing function}.
\end{definition}
=
\begin{proposition}[Interior differentiation]
On the interior $\mathcal{H}^{\circ}$, $A$ is infinitely differentiable, and
\begin{equation}
\nabla A(\eta)=\mathbb{E}_{\eta}T(X),
\qquad
\nabla^2 A(\eta)=\operatorname{Cov}_{\eta}\bigl(T(X)\bigr).
\tag{6}
\end{equation}
For $\eta$ in this interior and $t$ sufficiently close to zero,
\[
\mathbb{E}_{\eta}e^{t^\top T(X)}
=
\exp\bigl\{A(\eta+t)-A(\eta)\bigr\}.
\]
\end{proposition}

\begin{proof}
By the definition of the log-normalizing function,
\begin{align*}
\nabla A(\eta)
&=
\frac{\displaystyle\int T(x)e^{\eta^\top T(x)}h(x)\,\mathrm d\nu(x)}
{\displaystyle\int e^{\eta^\top T(x)}h(x)\,\mathrm d\nu(x)}\\
&=\int T(x)e^{\eta^\top T(x)-A(\eta)}h(x)\,\mathrm d\nu(x)\\
&=\mathbb E_\eta T(X).
\end{align*}
For the Hessian, differentiating the $i$th component of the gradient with
respect to $\eta_j$ gives
\begin{align*}
\bigl[\nabla^2 A(\eta)\bigr]_{ij}
&=\partial_j\int T_i(x)e^{\eta^\top T(x)-A(\eta)}h(x)\,\mathrm d\nu(x)\\
&=\int T_i(x)\bigl(T_j(x)-\partial_j A(\eta)\bigr)
e^{\eta^\top T(x)-A(\eta)}h(x)\,\mathrm d\nu(x)\\
&=\mathbb E_\eta[T_i(X)T_j(X)]
-\mathbb E_\eta[T_i(X)]\mathbb E_\eta[T_j(X)]\\
&=\operatorname{Cov}_\eta\bigl(T_i(X),T_j(X)\bigr).
\end{align*}
Thus
\[
\nabla^2 A(\eta)=\operatorname{Cov}_\eta\bigl(T(X)\bigr).
\]
Finally, for $t$ sufficiently close to zero so that $\eta+t\in\mathcal H$,
\begin{align*}
\mathbb E_\eta e^{t^\top T(X)}
&=\int e^{t^\top T(x)}e^{\eta^\top T(x)-A(\eta)}h(x)\,\mathrm d\nu(x)\\
&=e^{-A(\eta)}\int e^{(\eta+t)^\top T(x)}h(x)\,\mathrm d\nu(x)\\
&=e^{A(\eta+t)-A(\eta)}.
\end{align*}

It remains to justify the differentiations under the integral. Since
$\eta\in\mathcal H^\circ$, choose $\delta>0$ small enough that
$\eta+\delta s\in\mathcal H$ for every $s\in\{-1,1\}^m$.
For every integer $q\geq0$, there is a finite constant $C_{q,\delta}$ such
that, uniformly for $\lVert\eta'-\eta\rVert_\infty\leq\delta/2$,
\[
\lVert T(x)\rVert^q e^{(\eta')^\top T(x)}h(x)
\leq C_{q,\delta}\sum_{s\in\{-1,1\}^m}
e^{(\eta+\delta s)^\top T(x)}h(x).
\]
Every term on the right is integrable by the definition of $\mathcal H$.
Hence dominated convergence applies to derivatives of every order, so the
preceding calculations are valid and $A\in C^\infty(\mathcal H^\circ)$.
\end{proof}

\begin{definition}[Decision rule, loss, and risk]
A \emph{decision rule} is a measurable map
\[
\delta:\mathcal X\longrightarrow\mathcal A,
\]
where $\mathcal A$ is the action space. It turns data into an estimate, a
prediction, or a testing decision. A \emph{loss function}
$L(\theta,a)\geq0$ specifies the cost of taking action $a$ when the parameter
is $\theta$; more generally, write $L(P,a)$. The \emph{risk function}
averages this cost over data generated under the same parameter:
\begin{equation}
R(\theta,\delta)
=\mathbb E_\theta L\bigl(\theta,\delta(X)\bigr)
=\int L\bigl(\theta,\delta(x)\bigr)\,\mathrm dP_\theta(x).
\tag{7}
\end{equation}
The parameter is held fixed in this expectation. Different rules give
different functions of that parameter, not just one performance number.
Under squared-error loss for a real target $\psi(\theta)$, the risk is its
\emph{mean squared error} (MSE):
\[
R(\theta,\delta)
=\mathbb E_\theta\left[\bigl(\delta(X)-\psi(\theta)\bigr)^2\right].
\]
We use $\delta$ for a decision rule to distinguish it from a generic summary
$T$; a deterministic decision rule is itself a statistic.
\end{definition}

\begin{definition}[Bias and unbiasedness]
For i.i.d.\ real observations with mean $\mu$ and finite variance $\sigma^2$,
the sample mean estimates $\mu$ under squared loss and has risk
\[
R(\mu,\bar X)
=\mathbb E\bigl[(\bar X-\mu)^2\bigr]
=\frac{\sigma^2}{n}.
\]
For an
integrable estimator $\delta(X)$ of a target $\psi(P)$, its \emph{bias} is
\[
b_\delta(P)=\mathbb E_P\delta(X)-\psi(P).
\]
It is \emph{unbiased} if this equals zero for every $P$ in the model. For a
square-integrable estimator,
\begin{equation}
\mathbb E_P\bigl[(\delta(X)-\psi(P))^2\bigr]
=\operatorname{Var}_P\bigl(\delta(X)\bigr)+b_\delta(P)^2.
\tag{9}
\end{equation}
Subtract and add $\mathbb E_P\delta(X)$ inside the square; the cross term
has expectation zero. An unbiased estimator need not have the smallest MSE.
The sample variance $S^2$ is unbiased for $\sigma^2$, while the empirical
variance with denominator $n$ is biased.
\end{definition}

\begin{definition}[Hypothesis testing and error probabilities]
Suppose $\mathcal P_0$ and $\mathcal P_1$ are disjoint and the model under
consideration is $\mathcal P_0\cup\mathcal P_1$. We test
$H_0:P\in\mathcal P_0$ against $H_1:P\in\mathcal P_1$.
A nonrandomized test takes values in $\{0,1\}$, with $1$ meaning reject
$H_0$. Its \emph{rejection region} is
\[
C=\{x:\delta(x)=1\},
\]
and the test is the indicator $\delta(X)=1_C(X)$, not the set $C$ itself.
Under $0$--$1$ loss, a correct decision costs zero and an incorrect one
costs one:
\[
L(P,j)=
\begin{cases}
0, & P\in\mathcal P_j,\\
1, & P\notin\mathcal P_j,
\end{cases}
\qquad j\in\{0,1\},
\]
\[
R(P,\delta)=
\begin{cases}
P(X\in C), & P\in\mathcal P_0,\\
P(X\notin C), & P\in\mathcal P_1.
\end{cases}
\]
These are the \emph{type I error} probability under the null and the
\emph{type II error} probability under the alternative. Not rejecting a
null is a decision, not a proof that it is true.
\end{definition}

\begin{definition}[Sufficient statistic]
A statistic $T(X)$ is \emph{sufficient} for $\mathcal P=\{P_\theta\}$ if
there is a single probability kernel $K(t,B)$, not depending on $\theta$,
such that for every $\theta$ and every Borel set $B$,
\[
P_\theta\bigl(X\in B\mid T(X)\bigr)
=K\bigl(T(X),B\bigr)
\qquad P_\theta\text{-almost surely}.
\]
Thus the conditional distribution of the data given the statistic can be
chosen to be the same across the model. Conditional laws at summary values
having zero probability are interpreted through these versions, not by
dividing probabilities of zero events.
\end{definition}

\begin{theorem}
    If $T$ is a suffient statistic and $S$ is another statistic such that there is a measurable function $\phi$ such that $T = \phi(S)$, then $S$ is sufficient.
\end{theorem}
\begin{proof}
    There is a some $K(t,B)$ such that
    \[
    P_{\theta}(X \in A|T(X)) = K(T(X),A),\quad P_{\theta}\text{-a.s.}
    \]
    Then we know there exists a r.c.d. $L$ such that
    \[
    P_{K(t,\cdot)}(X \in A|S) = L(t,S,A)
    \]
    and hence
    \[
    K(t,A\cap S^{-1}(B)) = \int_{S^{-1}(B)} L(t,S(x),A)K(t,dx).
    \]
    We may say that
    \[
    K(T())
    \]
\end{proof}

\begin{theorem}[Neyman--Fisher factorization]
Let $\{P_\theta\}$ be a family on a Euclidean Borel sample space, dominated
by a common sigma-finite measure $\nu$. A Borel statistic $T$ with values in
a Euclidean space is sufficient if and only if its densities admit the
representation
\begin{equation}
p_\theta(x)=g_\theta\bigl(T(x)\bigr)h(x)
\qquad \nu\text{-almost everywhere for each }\theta,
\tag{10}
\end{equation}
where $h$ and $g_\theta$ are nonnegative measurable functions and $h$ does
not depend on $\theta$.
\end{theorem}

\newpage